\chapter{Gradually and indirectly, over time, gain in political significance (I'd Never Allow That to Happen)}
Magical girls have always been witches.

In the extradiegetic, historical sense, this is clearly true. The magical girl genre emerged as a direct result of the surprise popularity of Bewitched with a generation of Japanese schoolgirls \cite{ref59}. Samantha is the archetypal magical girl, conventionally attractive, traditionally feminine, with tremendous power tightly constrained within a limited sphere, and subject to the anxious masculinity that hedges her into that sphere. She can—and frequently does—assert herself, but ultimately she is trapped by the limitations of the norms of television in her time, bound to perform femininity with perfect makeup and hair, cute dresses, and a socially approved role as a wife and mother \cite{ref60}. 

While there are manga examples predating her, the first animated magical girl was likewise a witch, in Sally the Witch. Much, much younger than Bewitched's Samantha, Sally of course does not take on a role as wife or mother, but instead performs a child's femininity, being sweet and cute and, despite her power, fundamentally harmless \cite{ref61}. 

So it went. By Sailor Moon, the norms of the genre were largely set. Magical girls, like witches, gain their power from otherworldly sources, whether granted by the ruler of the Mirror Kingdom, accidentally released from a mysterious book, or inborn as a result of their past lives in the royal courts of the Moon. Magical girls, like witches, have their familiars, sentient creatures taking on animal form. And of course, magical girls, like witches, wield tremendous and varied magical powers.

But curiously, those powers are always directed at enemies as otherworldly as the magical girl's origins—indeed, often the enemies are tied intimately to those origins. As a general rule, magical girls do not fight government corruption, corporate malfeasance, or even that mainstay of the masculine hero, street crime. Their power, in other words, not only comes from the fantastic, but can only be directed against the fantastic.

Which is the other sense in which all magical girls are witches. The figure of the witch is a symbol of the fear of female power; in a world where masculinity is identified with hegemony and dominance, femininity must be identified with powerlessness, submission, or restraint \cite{ref62}. The traditional expression of this idea in Japanese culture is the figure of the ryosai kenbo, the "good wife, wise mother," a feminine ideal for which Hitomi appears to be training in Madoka Magica. Possessed of tremendous social intelligence, the ryosai kenbo rules utterly in the domestic sphere, having mastered many arts, but all are for the pleasure or support of her family, particularly supporting her husband and raising her children. Said husband, meanwhile, is the only one who is allowed to assert power outside of that sphere \cite{ref63}. Japanese folk and pop culture are rife with tales of the "bad" woman who wields power outside "her place," from the wife who turns out to be a shapeshifting kitsune \cite{ref64} to the cannibalistic old mountain hag \cite{ref65} to the seductive snow-woman who sucks the life-giving warmth from her paramour \cite{ref66}. 

Of course the good girl/bad girl dichotomy is hardly unique to Japanese culture. In Western folklore and pop culture it is represented (among a multitude of other representations) by the passive princess waiting to be rescued and the wicked witch who threatens her and the hero alike. To wield power is inherently to be the bad girl, the witch, a menace to the status quo. \cite{ref67} 

The power of the magical girl is usually sanitized in two ways. First, as already mentioned, her power is not permitted to impact anything the viewer might recognize as part of reality, but is instead almost invariably focused on fantastical opponents. Second, as we discussed back in Chapter 2, she is made to constantly perform femininity (remember, the counterpart to hegemonic masculinity is performed femininity), with frilly or skimpy costumes, elaborate poses, and of course the nude dance of the transformation sequence all serving to remind any potentially intimidated male viewers that she is still subject to the Male Gaze and still submitting to the social norms of the "good girl."

It is no accident that Sayaka's transformation into a witch immediately follows her using her powers on a pair of misogynists. She has stepped outside the boundaries of the good girl and challenged the status quo, and therefore is a bad girl, a wicked witch.

But this episode interrogates and ultimately subverts that binary in multiple ways. The most striking comes during Kyoko's fight with Oktavia, Sayaka's witch form; we see blue and red swirls of blood forming stylized images of Sayaka and Kyoko, which then swirl together into a rose, highly reminiscent of the opening to Revolutionary Girl Utena \cite{ref68}. That series also had a princess, the Rose Bride, who turned out to be a witch, and who (side by side with a swordwielding tomboy that positioned herself as the protector and rejected the usual feminine role) ultimately passed from submissive "good girl" to powerful and treacherous "bad girl" before finally breaking free of the entire system \cite{ref69}. That this is to be taken as a universal seems likely, given that said blood then splashes down in a shot framed to look like it is flowing from between Kyoko's legs.

In Western culture, the menstrual cycle has sometimes been posited as a particular punishment to women for their innate "badness," \cite{ref70} because of course it is the nature of performed femininity that to insist on being true to oneself is "bad" and leads inevitably to the label of witch. Thus all women have a "bad" streak, which is to say a coherent self that seeks expression.

But if magical girls and witches are truly one and have always been one, then what are we to make of Sayaka's transformation? Fortunately, the episode gives us the answer: it is the result of the system imposed by Kyubey. The magical girls' entire world has been imposed on them by the one significant male character, who holds total hegemony over them. His argument is that they have consented to take part in his system, which is of course absurd since he deliberately concealed crucial information from them; there are distinct shades of rape culture at work here, in the sense of the hegemonic male employing complex and nonsensical standards for what comprises consent, manipulating these definitions to place the blame on the victim.

But remember, Kyubey is a signifier of the implied author, who is not truly Gen Urobuchi but a gestalt entity formed from the combined efforts of writer, director, character designer, animators, voice actors, composer, and so on, an entire industry of creators. He spent this episode tricking Kyoko into treating Sayaka as someone to be saved, when he knows that Sayaka cannot be saved—after all, if she's been a witch all along, what is there to save her from? After successfully manipulating Kyoko into taking over the protector role—the same role which Sayaka was trying to fill—he tries to persuade Madoka to sacrifice herself similarly.

There is a term in anime fandom for a character (or, more accurately, a character trait) that invokes this protectiveness: moe. From the Japanese word "moeru," originally meaning "to burn" or "to scorch," it lent itself to meaning a passionate, burning feeling of protectiveness aroused by a helpless or endangered other. For an extended period in the late 2000s, an aesthetic rooted in that concept grew to dominate anime in general and magical girls in particular, though Saito locates its origins a decade earlier with otaku reactions to the Sailor Moon character Hotaru Tomoe (making the term something of a pun on her name). \cite{ref71} According to this aesthetic, the value of a character lies in their ability to evoke this protectiveness in a presumed-male audience, and the features which evoke it are helplessness, "cuteness," emotional vulnerability, and weakness, all coupled with a conventional and generic attractiveness \cite{ref72}. This is, of course, an extension of the same process that put Samantha under the thumb of her milquetoast husband and forces Sailor Moon to strip naked before she can access her powers; it renders the character harmless and therefore a "good girl," non-threatening to the inherent anxiety innate to hegemonic masculinity.

Kyubey stands revealed as the representative of a system that robs women of their power and makes them perform for his benefit, while also placing them into a position where their suffering is seen as proof of their need to be protected, which robs them of their power still further. He is, in other words, serving as an avatar of gender roles themselves. However, he is simply an instance of a larger and vaster system, which extends far beyond him; Madoka Magica itself is enmeshed in a larger culture, and while it can criticize ugly choices made in the name of economic viability, it cannot entirely escape them.

Nonetheless, the episode remains remarkably consistent. Sayaka's attempt to save others, to be the protector, transforms her into a monster. Kyoko and Madoka's attempt to save Sayaka gets them killed. Kyubey's quest to save the universe perpetuates a destructive and miserable system. And as we will see in the next episode, Homura's attempts to save Madoka are likewise doomed.

To seek to protect or save another, it seems, is inherently to rob them of power. But at the same time, the series has repeatedly vilified Kyubey for his lack of empathy, so it cannot be endorsing Objectivism, Ayn Rand's philosophy that includes rejecting altrusm as being immoral \cite{ref73}. Is there, perhaps, a difference between helping and saving? Or, as Urobuchi wrote in the Fate/zero author's notes, is it simply that we are helpless, and everything is doomed to become worse as all systems, universes, societies, and psyches alike, hurtle toward heat-death \cite{ref74}? 
